\section{Conclusion}
\label{sec:conclusion}

\subsection{Summary}

\fr{} is the first \recom{} variant that targets the uniform distribution over
valid redistricting plans as its stationary distribution.
It achieves this by augmenting each \recom{} step with a Metropolis-Hastings
correction: a second independent spanning tree sample from the proposed district
pair estimates the reverse transition probability, and the ratio of valid cut
counts determines the acceptance probability.

The key properties established in this paper:

\paragraph{Correct in expectation.}
The single-sample ratio $c_{\mathrm{fwd}} / c_{\mathrm{rev}}$ is an unbiased
estimator of the true MH ratio (Theorem~\ref{thm:edb}).
The chain satisfies detailed balance in expectation, targeting the uniform
distribution asymptotically.

\paragraph{Never fully stalls.}
By Proposition~\ref{prop:ratio-bounded}, the ratio is bounded away from 0
and $\infty$ whenever both spanning trees have at least one valid balanced cut.
The early-exit conditions (skip steps with zero valid cuts) prevent division
by zero and guarantee forward progress.

\paragraph{Verifiable on small graphs.}
The chi-squared test on the $2 \times 4$ grid confirms distributional
correctness: \fr{} is not rejected as targeting the uniform distribution
($p \approx 0.98$), while standard \recom{} is decisively rejected
($p \ll 10^{-6}$).

\paragraph{$2\times$ per-step cost.}
The Metropolis-Hastings correction requires one additional Wilson \ust{} call
per step.
Combined with the lower acceptance rate ($\approx\!40$--$60\%$ vs.
$\approx\!80\%$), \fr{} achieves approximately $28\%$ the accepted-step
throughput of standard \recom{} per unit wall-clock time.

\subsection{Usage Guidelines}

\begin{enumerate}
  \item \textbf{Use \fr{} when distributional correctness matters.}
        If the ensemble will be used to estimate the fraction of valid plans
        with a given partisan outcome, and the standard \recom{} bias is a
        concern, prefer \fr{}.
        CLI: \texttt{--search forest-recom}.

  \item \textbf{Use standard \recom{} for optimisation.}
        If the goal is to find good plans (low EC, high compactness) rather
        than to sample from the plan distribution, standard \recom{} or
        \sburst{}~\citep{dellaLibera2026shortburst} is faster and equally
        effective.

  \item \textbf{Use \smc{} for strict distributional guarantees.}
        For applications requiring exact calibration (e.g., formal expert
        testimony), SMC~\citep{dellaLibera2026smc} provides stronger guarantees
        at higher computational cost.
        \fr{} is a practical intermediate when SMC is infeasible.

  \item \textbf{Default step count.}
        We recommend at least $1{,}000$ steps for exploratory analysis and
        $5{,}000$--$10{,}000$ for distributional claims.
        For large states ($n > 4{,}000$, $k > 30$), the effective mixing time
        may require more steps; monitor convergence with $\hat{R}$ diagnostics.
\end{enumerate}

\subsection{Open Questions}

\begin{enumerate}
  \item \textbf{Variance of the ratio estimator.}
        The single-sample ratio $c_{\mathrm{fwd}} / c_{\mathrm{rev}}$ can have
        high variance when the true ratio is far from 1.
        Averaging over $m > 1$ independent reverse trees would reduce the
        variance at the cost of $m$ additional \ust{} calls per step.
        What is the optimal $m$ for typical census-tract graphs?

  \item \textbf{Mixing time under expected detailed balance.}
        Standard detailed balance guarantees mixing time bounds via spectral
        gap analysis.
        Expected detailed balance is weaker; what mixing time bounds follow
        from it?
        Are there gap results that apply to the \fr{} chain on planar graphs?

  \item \textbf{Non-uniform target distributions.}
        The \fr{} framework readily generalises to non-uniform targets $\mu(\pi)$
        by including the ratio $\mu(\pi') / \mu(\pi)$ in the acceptance probability.
        This would allow targeting distributions weighted by compactness or
        partisan fairness criteria.
        Implementing this generalisation in \texttt{redist-ensemble} is
        straightforward given the existing \fr{} infrastructure.

  \item \textbf{Forest ReCom and Short-Burst.}
        Running \sburst{}~\citep{dellaLibera2026shortburst} with \fr{} as the
        base chain (instead of standard \recom{}) would combine distributional
        correction with trajectory-level optimisation.
        This is the subject of G.12~\citep{dellaLibera2026shortburst}.
\end{enumerate}
